\section{Ablöseregelungen}\label{sec:abloseregelungen}

Aus vielen Gründen kann es vorkommen, dass die ganze Regierung (Legislative und Executive) oder auch nur Teile davon
vorzeitig abgelöst werden müssen.
Damit das Land zu keinem Zeitpunkt ohne Regierung bleibt, muss es für alle denkbaren Fälle Ablöseregelungen geben.
Die aktuell gültigen Regeln können nicht unmodifiziert übernommen werden.
Dazu sind die vorgeschlagenen Verfahrensänderungen zu grundlegend.
Zum Beispiel soll der Bundeskanzler jetzt ja indirekt vom Volk gewählt werden und nicht mehr vom Bundestag.
Daher kann der Bundestag nicht einfach so einen neuen Bundeskanzler wählen.
Das käme nur als befristete Übergangslösung in Betracht.

Im Folgenden werden daher die denkbaren Fälle einzeln durchgegangen:

\subsection{Ablösung des Bundeskanzlers ohne Legislativänderung}\label{subsec:ablosung-des-bundeskanzlers-ohne-legislativanderung}
- Kanzlerneuwahlen
\subsection{Zerbruch einer Koalition mit anschließender Bildung einer neuen Koalition}
Neuwahl des Bundeskanzlers durch den Bundestag nach Methode
\begin{enumerate}[label=\alph*)]
  \item Derjenige Kandidat der stärksten Koalitionspartei, der die meisten Stimmen erhalten hat, wird Bundeskanzler.
  \item Derjenige Kandidat von allen Koalitionsparteien, der die meisten Stimmen erhalten hat, wird Bundeskanzler.
  \item Der alte Bundeskanzler bleibt
\end{enumerate}
Wenn keine Wahl erfolgreich: Neuwahlen

\subsection{Zerbruch einer Koalition ohne Bildung einer neuen Koalition}
- Minderheitsregierung
- Neuwahlen

\subsection{Vorzeitige Ablösung einer Einparteienregierung}
- Neue Koalition
- Neuwahlen

\subsection{Ablösung einzelner Minister}
- Bundeskanzler entlässt und ernennt neu

\subsection{Ablösung einzelner Parlamentarier}
- Listenkandidaten rücken nach
- Direktwahlkandidaten: Personennachwahl im zugewiesenen Wahlkreis
